% bismi ALLAH ARRAHMAN ARRAHIM




















































\documentclass[11pt,a4paper]{article}

% =========================================================
% EX9 — Architecture de calcul modulaire à cores de sondes
% =========================================================

\usepackage[utf8]{inputenc}
\usepackage[T1]{fontenc}
\usepackage[french]{babel}
\usepackage{lmodern}
\usepackage{microtype}

\usepackage{geometry}
\geometry{
  left=2.3cm,
  right=2.3cm,
  top=2.2cm,
  bottom=2.2cm
}

\usepackage{amsmath,amssymb,mathtools}
\usepackage{booktabs}
\usepackage{array}
\usepackage{longtable}
\usepackage{tabularx}
\usepackage{enumitem}
\usepackage{xcolor}
\usepackage{graphicx}
\usepackage{float}
\usepackage{caption}
\usepackage{hyperref}

\usepackage{tikz}
\usetikzlibrary{
  arrows.meta,
  positioning,
  fit,
  calc,
  shapes.geometric,
  shapes.misc,
  backgrounds,
  matrix
}

\usepackage{pgfplots}
\pgfplotsset{compat=1.18}

\hypersetup{
  colorlinks=true,
  linkcolor=blue!50!black,
  citecolor=blue!50!black,
  urlcolor=blue!50!black,
  pdftitle={EX9 -- Architecture de calcul modulaire a cores de sondes},
  pdfauthor={Projet EX}
}

\setlength{\parindent}{0pt}
\setlength{\parskip}{5pt}
\setlist{nosep}

% Diagrammes TikZ : styles communs
\tikzset{
  probe/.style={
    draw=blue!65!black,
    fill=blue!7,
    rounded corners=2pt,
    minimum width=1.65cm,
    minimum height=0.68cm,
    align=center,
    font=\small
  },
  core/.style={
    draw=black!75,
    fill=gray!8,
    rounded corners=5pt,
    thick,
    inner sep=8pt
  },
  control/.style={
    draw=orange!70!black,
    fill=orange!10,
    rounded corners=3pt,
    minimum width=2.6cm,
    minimum height=0.8cm,
    align=center,
    font=\small
  },
  memory/.style={
    draw=green!50!black,
    fill=green!8,
    cylinder,
    shape border rotate=90,
    aspect=0.25,
    minimum width=1.5cm,
    minimum height=1.0cm,
    align=center,
    font=\small
  },
  data/.style={
    draw=purple!65!black,
    fill=purple!8,
    rounded corners=3pt,
    minimum width=1.8cm,
    minimum height=0.7cm,
    align=center,
    font=\small
  },
  flow/.style={
    -{Stealth[length=2mm]},
    thick
  },
  comm/.style={
    {Stealth[length=2mm]}-{Stealth[length=2mm]},
    dashed,
    thick
  }
}

\title{
  \textbf{EX9}\\[0.3cm]
  \Large Architecture de calcul modulaire\\
  \Large à cores composés de sondes spécialisées
}
\author{Projet EX}
\date{Document de conception et plan expérimental}

\begin{document}
\maketitle

\begin{abstract}
EX9 propose une architecture de calcul dans laquelle les unités
fonctionnelles ne sont plus considérées uniquement comme des couches
numériques monolithiques. Elles sont organisées en cores spécialisés,
chaque core étant lui-même composé de sondes.

L'architecture distingue trois variantes :
(i) les cores logiques, organisés dans un graphe logiciel ;
(ii) les cores physiques, constitués d'unités matérielles dédiées ;
(iii) les cores neuronaux, construits à partir de sondes adaptatives
et destinés à remplacer des blocs MLP.

Le calcul est représenté comme un graphe d'opérations spécialisées.
Des mécanismes de composition, de coordination, de stockage temporaire
et de communication organisent l'exécution de ce graphe.

L'objectif scientifique est de déterminer si cette organisation peut
réduire les calculs redondants, améliorer la spécialisation et offrir
un meilleur compromis entre précision, latence, mémoire et coût
d'exécution. Ces avantages constituent des hypothèses expérimentales,
et non des propriétés déjà démontrées.
\end{abstract}

\tableofcontents
\newpage

% =========================================================
\section{Positionnement et continuité avec EX7}
% =========================================================

\subsection{Point de départ expérimental}

EX7 V003, désigné \textit{OuterCoord-Lite-7x1}, constitue le point
de départ architectural disponible.

Cette version comprend :
\begin{itemize}
  \item sept sondes indépendantes de type \texttt{AdaptiveRBF} ;
  \item un coordinateur global unique ;
  \item une pondération globale apprise, notée $\rho$ ;
  \item une entrée commune transmise aux sept sondes ;
  \item une fusion des sorties des sondes.
\end{itemize}

Chaque sonde possède notamment :
\begin{itemize}
  \item un encodeur appris ;
  \item des centres RBF ;
  \item des largeurs adaptatives ;
  \item une couche de lecture produisant les logits.
\end{itemize}

Dans EX7 V003, les sept sondes sont toutes exécutées à l'inférence.
La coordination utilise une pondération globale fixe pour le modèle :
elle ne sélectionne pas dynamiquement les sondes selon chaque entrée.

EX7 démontre donc la faisabilité d'un ensemble de sondes adaptatives
coordonnées. Il ne démontre pas encore les gains d'une architecture
à cores spécialisés.

\subsection{Évolution proposée}

EX9 introduit une organisation hiérarchique :

\[
\boxed{
\text{Sondes}
\longrightarrow
\text{Cores}
\longrightarrow
\text{Système de cores}
}
\]

Un core est défini comme un ensemble structuré de sondes qui coopèrent
pour réaliser une fonction identifiable.

La notion de core ne désigne donc pas nécessairement un coeur CPU.
Elle désigne d'abord une unité fonctionnelle composée de sondes.
Selon la variante, cette unité peut être logicielle, physique ou
neuronale.

\subsection{Objectifs}

Les objectifs d'EX9 sont les suivants :

\begin{enumerate}
  \item Décomposer le calcul en opérations spécialisées.
  \item Organiser les opérations sous forme de graphes composables.
  \item Réutiliser les résultats intermédiaires.
  \item Séparer calcul, composition, coordination et stockage.
  \item Étudier l'exécution parallèle et conditionnelle.
  \item Construire des cores neuronaux capables de remplacer des
        blocs MLP.
  \item Étudier le passage d'une organisation logicielle à une
        organisation matérielle.
\end{enumerate}

% =========================================================
\section{Définition formelle d'une sonde}
% =========================================================

\subsection{Sonde comme opérateur}

Soit un espace d'entrée $\mathcal{X}$ et un espace de sortie
$\mathcal{Y}$.

Une sonde élémentaire est une application paramétrée :

\[
P_{\theta} :
\mathcal{X}
\rightarrow
\mathcal{Y}.
\]

Elle transforme une entrée en sortie selon des paramètres
$\theta$.

Dans le cas d'une sonde spécialisée dans une opération numérique,
on peut écrire :

\[
P_{\theta}^{(k)}(x)
\simeq
\phi_k(x),
\]

où $\phi_k$ représente l'opération visée.

Exemples :

\[
\begin{aligned}
P^{(\mathrm{PROD})}(a,b) &\simeq ab,\\
P^{(\mathrm{EXP})}(c) &\simeq e^c,\\
P^{(\mathrm{DIV})}(u,v) &\simeq \frac{u}{v},\\
P^{(\mathrm{LOG})}(j) &\simeq \log(j).
\end{aligned}
\]

Le symbole $\simeq$ indique une approximation éventuelle lorsque
la sonde est apprise. Pour une unité arithmétique exacte, on pourra
au contraire imposer une sémantique déterministe.

\subsection{Sonde adaptative de type RBF}

Une sonde neuronale peut utiliser une représentation latente :

\[
z = \tanh(Wx+b),
\]

où :
\[
W\in\mathbb{R}^{d\times n},
\qquad
b\in\mathbb{R}^{d}.
\]

Pour $K$ centres $c_k$, les activations radiales sont :

\[
r_k(z)
=
\exp\left(
-\frac{\|z-c_k\|_2^2}{2\sigma_k^2}
\right).
\]

La sortie de la sonde peut être définie par :

\[
P_{\theta}(x)
=
R(x)V+\beta,
\]

avec :

\[
R(x)=[r_1(z),\ldots,r_K(z)].
\]

Les paramètres appris comprennent notamment :

\[
\theta=(W,b,C,\sigma,V,\beta).
\]

Cette formulation reprend le principe de la sonde
\texttt{AdaptiveRBF} utilisée dans EX7.

Elle ne signifie pas que toutes les sondes EX9 devront être des RBF.
Les sondes pourront être numériques exactes, neuronales adaptatives,
ou spécialisées dans le routage et la gestion de données.

% =========================================================
\section{Définition formelle d'un core}
% =========================================================

\subsection{Définition générale}

Un core $C_i$ est un ensemble structuré de $m_i$ sondes :

\[
C_i =
\left(
P_{i,1},
P_{i,2},
\ldots,
P_{i,m_i},
\Gamma_i,
\mathcal{M}_i
\right),
\]

où :

\begin{itemize}
  \item $P_{i,j}$ est une sonde appartenant au core ;
  \item $\Gamma_i$ décrit l'organisation et les dépendances internes ;
  \item $\mathcal{M}_i$ décrit les mécanismes de mémoire ou d'état
        associés au core.
\end{itemize}

Le core implémente une fonction globale :

\[
C_i : \mathcal{X}_i \rightarrow \mathcal{Y}_i.
\]

Cette fonction résulte de la composition des sondes et de leurs
interactions.

Un core peut être spécialisé dans :
\begin{itemize}
  \item le calcul numérique ;
  \item la composition des opérations ;
  \item la coordination ;
  \item le stockage temporaire ;
  \item la communication ;
  \item le traitement neuronal adaptatif.
\end{itemize}

\subsection{Structure interne d'un core}

\begin{figure}[H]
\centering
\resizebox{0.92\textwidth}{!}{
\begin{tikzpicture}[node distance=0.45cm]
  \node[probe] (p1) {Sonde 1};
  \node[probe,right=of p1] (p2) {Sonde 2};
  \node[probe,right=of p2] (p3) {Sonde 3};
  \node[probe,right=of p3] (p4) {Sonde $m$};

  \node[core,fit=(p1)(p2)(p3)(p4),
        label=above:{\textbf{CORE}}] (c) {};

  \node[data,below=0.8cm of c] (in) {Entrées};
  \node[data,right=3.0cm of in] (out) {Sorties};

  \draw[flow] (in) -- (p1.south);
  \draw[flow] (p4.south) -- (out);
  \draw[comm] (p1) -- (p2);
  \draw[comm] (p2) -- (p3);
  \draw[comm] (p3) -- (p4);
\end{tikzpicture}
}
\caption{Un core est composé de sondes qui coopèrent.}
\end{figure}

Cette représentation est conceptuelle : les sondes ne sont pas
obligatoirement exécutées en série. Leur organisation peut être
parallèle, hiérarchique ou conditionnelle.

% =========================================================
\section{Les trois variantes de cores}
% =========================================================

\subsection{Core logique}

Un core logique est une unité fonctionnelle logicielle constituée
de sondes et de mécanismes de composition.

Il n'impose aucune correspondance directe avec un coeur matériel.
Plusieurs cores logiques peuvent être exécutés sur un même CPU ou GPU.

Formellement :

\[
C_i^{(L)}
=
(\mathcal{P}_i,\Gamma_i,\mathcal{S}_i),
\]

où :
\begin{itemize}
  \item $\mathcal{P}_i$ est l'ensemble des sondes ;
  \item $\Gamma_i$ est le graphe d'exécution ;
  \item $\mathcal{S}_i$ est l'état logiciel du core.
\end{itemize}

Le graphe d'exécution est :

\[
G_i=(V_i,E_i),
\]

avec :
\begin{itemize}
  \item $V_i$ : sondes et opérations ;
  \item $E_i$ : dépendances de données et de contrôle.
\end{itemize}

\begin{figure}[H]
\centering
\resizebox{0.88\textwidth}{!}{
\begin{tikzpicture}[node distance=0.65cm]
  \node[data] (x) {Entrée};
  \node[probe,right=of x] (p1) {Sonde A};
  \node[probe,above right=0.55cm and 0.8cm of p1] (p2) {Sonde B};
  \node[probe,below right=0.55cm and 0.8cm of p1] (p3) {Sonde C};
  \node[probe,right=1.0cm of p2] (p4) {Sonde D};
  \node[data,right=1.0cm of p4] (y) {Sortie};

  \draw[flow] (x) -- (p1);
  \draw[flow] (p1) -- (p2);
  \draw[flow] (p1) -- (p3);
  \draw[flow] (p2) -- (p4);
  \draw[flow] (p3) -- (p4);
  \draw[flow] (p4) -- (y);

  \begin{scope}[on background layer]
    \node[core,fit=(p1)(p2)(p3)(p4),
          label=above:{Core logique}] {};
  \end{scope}
\end{tikzpicture}
}
\caption{Exemple d'un core logique organisé en graphe logiciel.}
\end{figure}

\textbf{Propriété :} le core logique définit l'organisation fonctionnelle,
mais ne garantit pas à lui seul une accélération matérielle.

\subsection{Core physique}

Un core physique correspond à une réalisation matérielle d'un ensemble
de sondes.

Il peut comprendre :
\begin{itemize}
  \item des unités de calcul ;
  \item des unités de composition ;
  \item une mémoire locale ;
  \item des mécanismes de communication ;
  \item un contrôleur d'exécution.
\end{itemize}

On définit :

\[
C_i^{(P)}
=
(\mathcal{H}_i,\mathcal{P}_i,\mathcal{M}_i,\mathcal{N}_i),
\]

où :
\begin{itemize}
  \item $\mathcal{H}_i$ désigne les ressources matérielles ;
  \item $\mathcal{P}_i$ désigne les sondes exécutées ;
  \item $\mathcal{M}_i$ désigne la mémoire locale ;
  \item $\mathcal{N}_i$ désigne le réseau de communication.
\end{itemize}

\begin{figure}[H]
\centering
\resizebox{0.92\textwidth}{!}{
\begin{tikzpicture}[node distance=0.5cm]
  \node[probe] (p1) {Unité sonde 1};
  \node[probe,right=of p1] (p2) {Unité sonde 2};
  \node[probe,right=of p2] (p3) {Unité sonde 3};

  \node[control,below=0.9cm of p2] (ctrl) {Contrôleur};
  \node[memory,below=0.8cm of p1] (mem) {Mémoire locale};
  \node[data,below=0.8cm of p3] (net) {Interface réseau};

  \node[core,fit=(p1)(p2)(p3)(ctrl)(mem)(net),
        label=above:{\textbf{CORE PHYSIQUE}}] (c) {};

  \draw[comm] (p1) -- (p2);
  \draw[comm] (p2) -- (p3);
  \draw[flow] (ctrl) -- (p2);
  \draw[comm] (mem) -- (p1);
  \draw[comm] (net) -- (p3);
\end{tikzpicture}
}
\caption{Organisation conceptuelle d'un core physique.}
\end{figure}

Le passage à cette variante nécessite une étude de :
\begin{itemize}
  \item la surface matérielle ;
  \item la consommation énergétique ;
  \item la bande passante mémoire ;
  \item la latence des communications ;
  \item le parallélisme réellement exploitable.
\end{itemize}

Un core physique n'est pas nécessairement plus rapide qu'un CPU ou GPU.
Le bénéfice doit être mesuré sur une réalisation concrète.

\subsection{Core neuronal}

Un core neuronal est un ensemble de sondes adaptatives qui réalise
une transformation apprise.

Il est conçu pour remplacer un bloc MLP ou une partie d'un réseau
neuronal classique.

On définit :

\[
C_i^{(N)}(x)
=
\mathcal{F}_{\Theta_i}
(P_{i,1},\ldots,P_{i,m_i};x),
\]

où $\mathcal{F}_{\Theta_i}$ est l'opérateur de composition neuronal.

La fonction du core est apprise à partir de données :

\[
\Theta_i^\star
=
\arg\min_{\Theta_i}
\mathbb{E}_{(x,y)\sim\mathcal{D}}
\left[
\mathcal{L}
\left(
C_i^{(N)}(x),y
\right)
\right].
\]

\begin{figure}[H]
\centering
\resizebox{0.90\textwidth}{!}{
\begin{tikzpicture}[node distance=0.55cm]
  \node[data] (x) {Vecteur $x$};

  \node[probe,right=of x,yshift=0.9cm] (p1) {Sonde adaptative 1};
  \node[probe,right=of x] (p2) {Sonde adaptative 2};
  \node[probe,right=of x,yshift=-0.9cm] (p3) {Sonde adaptative 3};

  \node[control,right=1.0cm of p2] (fusion) {Composition\\neurale};
  \node[data,right=of fusion] (y) {Sortie};

  \draw[flow] (x) -- (p1);
  \draw[flow] (x) -- (p2);
  \draw[flow] (x) -- (p3);
  \draw[flow] (p1) -- (fusion);
  \draw[flow] (p2) -- (fusion);
  \draw[flow] (p3) -- (fusion);
  \draw[flow] (fusion) -- (y);

  \begin{scope}[on background layer]
    \node[core,fit=(p1)(p2)(p3)(fusion),
          label=above:{\textbf{CORE NEURONAL}}] {};
  \end{scope}
\end{tikzpicture}
}
\caption{Core neuronal composé de sondes adaptatives.}
\end{figure}

\textbf{Distinction importante :} une sonde n'est pas nécessairement
un neurone scalaire. Une sonde adaptative peut être un petit module
vectoriel possédant son propre encodeur, ses paramètres et sa fonction
de sortie.

\subsection{Comparaison des trois variantes}

\begin{table}[H]
\centering
\small
\begin{tabularx}{\textwidth}{p{2.5cm}XXX}
\toprule
 & Core logique & Core physique & Core neuronal \\
\midrule
Nature
& Logicielle
& Matérielle
& Apprise \\

Composition
& Sondes et graphe
& Sondes et unités matérielles
& Sondes adaptatives \\

Exécution
& CPU/GPU
& Unités physiques dédiées
& Calcul neuronal \\

Objectif
& Organiser
& Accélérer matériellement
& Remplacer des blocs MLP \\

Validation
& Coût logiciel
& Latence, énergie, surface
& Qualité, paramètres, latence \\
\bottomrule
\end{tabularx}
\caption{Différences entre les variantes de cores EX9.}
\end{table}

Ces trois variantes sont complémentaires et peuvent être combinées.
Un core logique peut être le modèle de référence d'un core physique.
Un core neuronal peut être exécuté dans un core logique ou physique.

% =========================================================
\section{Calcul sous forme de graphe de sondes}
% =========================================================

\subsection{Calcul classique}

Considérons :

\[
f(a,b,c,j)
=
ab+\frac{e^c}{a}+\log(j).
\]

Dans une expression classique, le calcul est décrit par une suite
d'opérations arithmétiques.

EX9 propose de représenter chaque opération par une sonde ou une
unité spécialisée.

\subsection{Graphe de calcul}

Définissons un graphe orienté acyclique :

\[
G=(V,E).
\]

Chaque sommet $v_i\in V$ représente une opération ou une sonde.
Chaque arête $(v_i,v_j)\in E$ indique qu'une sortie de $v_i$ est
utilisée par $v_j$.

Pour l'expression précédente :

\[
\begin{aligned}
u_1 &= P_{\mathrm{PROD}}(a,b),\\
u_2 &= P_{\mathrm{EXP}}(c),\\
u_3 &= P_{\mathrm{DIV}}(u_2,a),\\
u_4 &= P_{\mathrm{LOG}}(j),\\
y &= P_{\mathrm{SUM}}(u_1,u_3,u_4).
\end{aligned}
\]

Si les sondes sont exactes, cette composition réalise l'expression
arithmétique initiale. Si elles sont apprises, une erreur
d'approximation peut apparaître à chaque étape.

\begin{figure}[H]
\centering
\resizebox{0.94\textwidth}{!}{
\begin{tikzpicture}[node distance=0.85cm and 0.75cm]
  \node[data] (a) {$a,b$};
  \node[data,below=0.65cm of a] (c) {$c$};
  \node[data,below=0.65cm of c] (j) {$j$};

  \node[probe,right=of a] (prod) {PROD};
  \node[probe,right=of c] (exp) {EXP};
  \node[probe,right=of j] (log) {LOG};

  \node[probe,right=1.2cm of exp] (div) {DIV};
  \node[probe,right=1.2cm of div] (sum) {SUM};

  \node[data,right=of sum] (y) {$y$};

  \draw[flow] (a) -- (prod);
  \draw[flow] (c) -- (exp);
  \draw[flow] (exp) -- (div);
  \draw[flow] (a.east) -- ++(0.45,0) |- (div.north);
  \draw[flow] (j) -- (log);
  \draw[flow] (prod) -- (sum);
  \draw[flow] (div) -- (sum);
  \draw[flow] (log) -- ++(0.45,0) |- (sum.south);
  \draw[flow] (sum) -- (y);
\end{tikzpicture}
}
\caption{Graphe de calcul composé de sondes spécialisées.}
\end{figure}

\subsection{Dépendances et parallélisme}

Les opérations PROD, EXP et LOG sont indépendantes au départ.
Elles peuvent donc être exécutées en parallèle.

L'opération DIV dépend de EXP et de la valeur $a$.
L'opération SUM dépend des trois branches.

Le graphe exprime donc les dépendances réelles du calcul, plutôt
qu'une séquence arbitraire.

On note $\operatorname{Pred}(v)$ l'ensemble des prédécesseurs d'un
sommet $v$. Une opération est exécutable lorsque toutes ses entrées
sont disponibles :

\[
\operatorname{Ready}(v,t)
=
\bigwedge_{u\in\operatorname{Pred}(v)}
\operatorname{Available}(u,t).
\]

Cette règle rend possible une exécution parallèle des opérations
indépendantes.

% =========================================================
\section{Composition des sondes}
% =========================================================

\subsection{Composition fonctionnelle}

Soient deux sondes :

\[
P_1:\mathcal{X}\rightarrow\mathcal{Z},
\qquad
P_2:\mathcal{Z}\rightarrow\mathcal{Y}.
\]

Leur composition est :

\[
P_{2\circ1}(x)=P_2(P_1(x)).
\]

Pour un graphe de calcul général, chaque sonde reçoit les sorties
de ses prédécesseurs :

\[
h_i
=
P_i
\left(
\{h_j:j\in\operatorname{Pred}(i)\}
\right).
\]

La sortie globale est produite par un ou plusieurs sommets terminaux.

\subsection{Composition hiérarchique}

Les sondes peuvent être regroupées en cores, puis les cores peuvent
être composés :

\[
\mathcal{P}
\longrightarrow
\{C_1,\ldots,C_M\}
\longrightarrow
F_{\mathrm{EX9}}.
\]

La fonction globale devient :

\[
F_{\mathrm{EX9}}(x)
=
\mathcal{C}
\left(
C_1(x_1),\ldots,C_M(x_M)
\right),
\]

où $\mathcal{C}$ est l'opérateur de composition inter-cores.

Les entrées $x_i$ peuvent être identiques, projetées, partitionnées
ou partagées selon la topologie retenue.

% =========================================================
\section{Réutilisation et stockage temporaire}
% =========================================================

\subsection{Problème de redondance}

Un même sous-calcul peut être utilisé plusieurs fois.

Par exemple :

\[
u=\exp(c)
\]

peut alimenter plusieurs opérations.

Une architecture sans mécanisme de réutilisation peut recalculer
plusieurs fois la même valeur.

EX9 propose un stockage temporaire des résultats intermédiaires.

\subsection{Mémoire temporaire}

On définit une mémoire temporaire :

\[
\mathcal{M}_t
=
\{(k,v,\tau,\ell)\},
\]

où :
\begin{itemize}
  \item $k$ est une clé identifiant le calcul ;
  \item $v$ est le résultat calculé ;
  \item $\tau$ est l'instant ou la version de calcul ;
  \item $\ell$ est la durée de validité ou la politique de conservation.
\end{itemize}

Une opération peut consulter la mémoire avant de lancer un calcul :

\[
\operatorname{Execute}(k,x)
=
\begin{cases}
\mathcal{M}_t[k],
& \text{si une valeur valide existe},\\
P_k(x),
& \text{sinon}.
\end{cases}
\]

Après calcul, le résultat peut être enregistré :

\[
\mathcal{M}_{t+1}
=
\mathcal{M}_t
\cup
\{(k,P_k(x),t,\ell)\}.
\]

Cette stratégie correspond à une forme de mémorisation des
sous-calculs, souvent appelée \textit{memoization}.

\subsection{Exemple de réutilisation}

\[
\begin{aligned}
u &= \operatorname{EXP}(c),\\
v_1 &= \operatorname{DIV}(u,a),\\
v_2 &= \operatorname{PROD}(u,b),\\
v_3 &= \operatorname{SUM}(v_1,v_2).
\end{aligned}
\]

La valeur $u$ est calculée une fois et réutilisée deux fois.

\begin{figure}[H]
\centering
\resizebox{0.90\textwidth}{!}{
\begin{tikzpicture}[node distance=0.8cm and 0.8cm]
  \node[data] (c) {$c$};
  \node[probe,right=of c] (exp) {EXP};
  \node[memory,right=of exp] (mem) {Mémoire\\temporaire};

  \node[probe,right=1.5cm of mem,yshift=0.8cm] (div) {DIV};
  \node[probe,right=1.5cm of mem,yshift=-0.8cm] (prod) {PROD};

  \node[probe,right=1.1cm of div,yshift=-0.8cm] (sum) {SUM};

  \draw[flow] (c) -- (exp);
  \draw[flow] (exp) -- (mem);
  \draw[comm] (mem) -- (div);
  \draw[comm] (mem) -- (prod);
  \draw[flow] (div) -- (sum);
  \draw[flow] (prod) -- (sum);
\end{tikzpicture}
}
\caption{Une valeur intermédiaire est calculée une fois puis réutilisée.}
\end{figure}

\subsection{Conditions de validité}

La réutilisation n'est correcte que si les entrées et le contexte
du calcul sont compatibles.

La clé de mémoire doit prendre en compte, selon le cas :
\begin{itemize}
  \item les valeurs d'entrée ;
  \item les paramètres de la sonde ;
  \item la version du modèle ;
  \item le contexte numérique ;
  \item les éventuels effets de bord.
\end{itemize}

Pour des sondes stochastiques ou dépendantes d'un état mutable,
une réutilisation naïve peut modifier le comportement du système.

La mémoire temporaire doit donc être accompagnée d'une politique
explicite de validité et d'invalidation.

% =========================================================
\section{Core de coordination}
% =========================================================

\subsection{Rôle}

Le core de coordination organise le système de cores.

Il peut :
\begin{itemize}
  \item sélectionner les cores nécessaires ;
  \item déterminer l'ordre d'exécution ;
  \item identifier les dépendances ;
  \item allouer les ressources ;
  \item décider du partage des résultats ;
  \item contrôler les transferts de données.
\end{itemize}

On représente la décision de coordination par :

\[
\pi_{\mathrm{coord}}:
(s_t,x_t,G)
\rightarrow
a_t,
\]

où :
\begin{itemize}
  \item $s_t$ est l'état courant du système ;
  \item $x_t$ est l'entrée ou la donnée disponible ;
  \item $G$ est le graphe de calcul ;
  \item $a_t$ est une action de coordination.
\end{itemize}

Une action peut être :

\[
a_t\in
\{
\text{ACTIVER},
\text{ATTENDRE},
\text{TRANSMETTRE},
\text{RÉUTILISER},
\text{STOCKER},
\text{FUSIONNER}
\}.
\]

\subsection{Architecture inter-cores}

\begin{figure}[H]
\centering
\resizebox{0.94\textwidth}{!}{
\begin{tikzpicture}[node distance=0.75cm and 0.9cm]
  \node[control] (coord) {Core de\\coordination};

  \node[core,below left=1.0cm and 2.1cm of coord] (calc) {
    \begin{minipage}{3.2cm}
      \centering
      \textbf{Core de calcul}\\[2mm]
      Sonde PROD\\
      Sonde EXP\\
      Sonde DIV\\
      Sonde LOG
    \end{minipage}
  };

  \node[core,below=1.0cm of coord] (comp) {
    \begin{minipage}{3.0cm}
      \centering
      \textbf{Core de composition}\\[2mm]
      Fusion\\
      Assemblage\\
      Dépendances
    \end{minipage}
  };

  \node[core,below right=1.0cm and 2.1cm of coord] (mem) {
    \begin{minipage}{3.0cm}
      \centering
      \textbf{Core mémoire}\\[2mm]
      Résultats\\
      Intermédiaires\\
      Réutilisation
    \end{minipage}
  };

  \node[core,below=1.1cm of comp] (comm) {
    \begin{minipage}{3.0cm}
      \centering
      \textbf{Core communication}\\[2mm]
      Routage\\
      Transfert\\
      Synchronisation
    \end{minipage}
  };

  \draw[comm] (coord) -- (calc);
  \draw[comm] (coord) -- (comp);
  \draw[comm] (coord) -- (mem);
  \draw[comm] (calc) -- (comp);
  \draw[comm] (comp) -- (mem);
  \draw[comm] (calc) -- (comm);
  \draw[comm] (mem) -- (comm);
  \draw[comm] (comp) -- (comm);
\end{tikzpicture}
}
\caption{Organisation fonctionnelle proposée pour EX9.}
\end{figure}

Les cores présentés sont des rôles fonctionnels. Ils peuvent être
regroupés, distribués ou fusionnés selon le coût réel de leur
implémentation.

Il n'est pas nécessaire d'imposer un core physique distinct pour
chaque fonction.

% =========================================================
\section{Langage de coordination exécutable}
% =========================================================

\subsection{Principe}

Le core de coordination peut manipuler un langage d'actions
décrivant les opérations à effectuer.

On définit un ensemble d'instructions :

\[
\mathcal{A}
=
\{
\operatorname{APPLY},
\operatorname{ROUTE},
\operatorname{STORE},
\operatorname{LOAD},
\operatorname{MERGE},
\operatorname{WAIT}
\}.
\]

Une instruction est interprétée par un exécuteur :

\[
\operatorname{Exec}:
(\mathcal{A},s_t)
\rightarrow
s_{t+1}.
\]

Une séquence d'actions :

\[
A=(a_1,a_2,\ldots,a_T)
\]

produit une évolution :

\[
s_{t+1}
=
\operatorname{Exec}(a_t,s_t).
\]

Le résultat final est :

\[
y=H(s_T).
\]

\subsection{Exemple de programme d'actions}

Pour calculer :

\[
y=ab+\frac{e^c}{a}+\log(j),
\]

le coordinateur peut produire une séquence conceptuelle :

\[
\begin{aligned}
a_1 &= \operatorname{APPLY}(\mathrm{PROD},a,b),\\
a_2 &= \operatorname{APPLY}(\mathrm{EXP},c),\\
a_3 &= \operatorname{STORE}(u),\\
a_4 &= \operatorname{APPLY}(\mathrm{DIV},u,a),\\
a_5 &= \operatorname{APPLY}(\mathrm{LOG},j),\\
a_6 &= \operatorname{MERGE}(u_1,u_2,u_3).
\end{aligned}
\]

Cette représentation ne supprime pas le calcul arithmétique.
Elle décrit comment les opérations sont sélectionnées, ordonnées,
routées et réutilisées.

\subsection{Condition d'efficacité}

Le langage d'actions n'est utile que si son coût total est inférieur
au coût de la méthode qu'il remplace.

On définit :

\[
T_{\mathrm{EX9}}
=
T_{\mathrm{selection}}
+
T_{\mathrm{coordination}}
+
T_{\mathrm{communication}}
+
T_{\mathrm{calcul}}
+
T_{\mathrm{mémoire}}.
\]

Le système n'est plus rapide que la référence si :

\[
T_{\mathrm{EX9}} < T_{\mathrm{baseline}}.
\]

Cette inégalité doit être vérifiée expérimentalement.

% =========================================================
\section{Architecture neuronale : remplacement des MLP}
% =========================================================

\subsection{MLP de référence}

Un bloc MLP classique peut être représenté par :

\[
\operatorname{MLP}(x)
=
W_2\sigma(W_1x+b_1)+b_2.
\]

Il effectue une transformation dense suivie d'une non-linéarité.

Dans EX9, un core neuronal doit apprendre à réaliser une fonction
comparable :

\[
C^{(N)}(x)
\approx
\operatorname{MLP}(x).
\]

\subsection{Core neuronal de remplacement}

Le core neuronal est composé de sondes adaptatives :

\[
C^{(N)}(x)
=
\mathcal{F}
\left(
P_1(x),P_2(x),\ldots,P_m(x)
\right).
\]

La fonction $\mathcal{F}$ peut être :
\begin{itemize}
  \item une somme pondérée ;
  \item une concaténation suivie d'une projection ;
  \item une composition hiérarchique ;
  \item un mécanisme de coordination appris ;
  \item une fusion conditionnelle.
\end{itemize}

L'objectif est d'obtenir une approximation fonctionnelle du bloc
MLP, sans supposer que la nouvelle architecture sera plus efficace.

\subsection{Comparaison fonctionnelle}

\begin{figure}[H]
\centering
\resizebox{0.94\textwidth}{!}{
\begin{tikzpicture}[node distance=0.65cm and 0.9cm]
  \node[data] (x1) {$x$};
  \node[probe,right=of x1,minimum width=2.5cm] (l1) {Projection dense};
  \node[probe,right=of l1,minimum width=2.5cm] (act) {Non-linéarité};
  \node[probe,right=of act,minimum width=2.5cm] (l2) {Projection dense};
  \node[data,right=of l2] (y1) {$y$};

  \draw[flow] (x1) -- (l1);
  \draw[flow] (l1) -- (act);
  \draw[flow] (act) -- (l2);
  \draw[flow] (l2) -- (y1);

  \node[below=0.35cm of act] {\textbf{MLP de référence}};
\end{tikzpicture}
}
\vspace{0.5cm}
\resizebox{0.94\textwidth}{!}{
\begin{tikzpicture}[node distance=0.55cm and 0.65cm]
  \node[data] (x) {$x$};
  \node[probe,right=of x,yshift=0.8cm] (p1) {Sonde 1};
  \node[probe,right=of x] (p2) {Sonde 2};
  \node[probe,right=of x,yshift=-0.8cm] (p3) {Sonde $m$};
  \node[control,right=1.1cm of p2] (f) {Composition\\adaptative};
  \node[data,right=of f] (y) {$y'$};

  \draw[flow] (x) -- (p1);
  \draw[flow] (x) -- (p2);
  \draw[flow] (x) -- (p3);
  \draw[flow] (p1) -- (f);
  \draw[flow] (p2) -- (f);
  \draw[flow] (p3) -- (f);
  \draw[flow] (f) -- (y);

  \node[below=0.35cm of f] {\textbf{Core neuronal EX9}};
\end{tikzpicture}
}
\caption{Remplacement expérimental d'un bloc MLP par un core neuronal.}
\end{figure}

\subsection{Protocole d'apprentissage}

Deux approches seront comparées.

\textbf{Apprentissage supervisé direct :}

\[
\Theta^\star
=
\arg\min_{\Theta}
\mathbb{E}_{(x,y)\sim\mathcal{D}}
\left[
\mathcal{L}(C_{\Theta}(x),y)
\right].
\]

\textbf{Distillation d'un MLP enseignant :}

\[
\Theta^\star
=
\arg\min_{\Theta}
\mathbb{E}_{x\sim\mathcal{D}}
\left[
D\left(
C_{\Theta}(x),
\operatorname{MLP}_{\mathrm{teacher}}(x)
\right)
\right].
\]

La distillation permet de comparer plus directement les fonctions
réalisées, indépendamment des différences d'optimisation.

% =========================================================
\section{Architecture globale EX9}
% =========================================================

\subsection{Vue d'ensemble}

L'architecture globale rassemble :
\begin{itemize}
  \item un core de coordination ;
  \item un ou plusieurs cores de calcul ;
  \item un core de composition ;
  \item un core de mémoire temporaire ;
  \item un core de communication ;
  \item éventuellement des cores neuronaux.
\end{itemize}

\begin{figure}[H]
\centering
\resizebox{0.98\textwidth}{!}{
\begin{tikzpicture}[node distance=0.65cm and 0.7cm]
  \node[data] (input) {Entrée};

  \node[control,right=of input] (coord) {Core de\\coordination};

  \node[core,right=1.0cm of coord,yshift=1.6cm] (calc) {
    \begin{minipage}{2.8cm}
      \centering
      \textbf{Core calcul}\\
      PROD / EXP / DIV / LOG
    \end{minipage}
  };

  \node[core,right=1.0cm of coord] (neural) {
    \begin{minipage}{2.8cm}
      \centering
      \textbf{Core neuronal}\\
      Sondes adaptatives
    \end{minipage}
  };

  \node[core,right=1.0cm of coord,yshift=-1.6cm] (comp) {
    \begin{minipage}{2.8cm}
      \centering
      \textbf{Core composition}\\
      Fusion et assemblage
    \end{minipage}
  };

  \node[memory,right=1.0cm of calc] (mem) {Mémoire\\temporaire};
  \node[data,right=1.0cm of mem] (output) {Sortie};

  \draw[flow] (input) -- (coord);
  \draw[comm] (coord) -- (calc);
  \draw[comm] (coord) -- (neural);
  \draw[comm] (coord) -- (comp);
  \draw[comm] (calc) -- (mem);
  \draw[comm] (neural) -- (mem);
  \draw[comm] (comp) -- (mem);
  \draw[flow] (mem) -- (output);
}
\end{tikzpicture}
}
\caption{Architecture globale conceptuelle EX9.}
\end{figure}

Cette figure représente les fonctions principales. Elle ne prescrit
pas une topologie matérielle unique.

\subsection{Formalisation globale}

Soit un ensemble de cores :

\[
\mathcal{C}
=
\{C_1,C_2,\ldots,C_M\}.
\]

Le système EX9 est défini par :

\[
\operatorname{EX9}
=
(\mathcal{C},G,\Pi,\mathcal{M},\mathcal{N}),
\]

où :
\begin{itemize}
  \item $\mathcal{C}$ est l'ensemble des cores ;
  \item $G$ est le graphe global de dépendances ;
  \item $\Pi$ est la politique de coordination ;
  \item $\mathcal{M}$ est l'organisation mémoire ;
  \item $\mathcal{N}$ est le mécanisme de communication.
\end{itemize}

Le système calcule :

\[
y
=
\operatorname{Execute}
\left(
x;\operatorname{EX9}
\right).
\]

La fonction $\operatorname{Execute}$ doit respecter les dépendances
du graphe, les règles de validité mémoire et les contrats d'entrée
et de sortie des cores.

% =========================================================
\section{Modèle de coût et hypothèses d'efficacité}
% =========================================================

\subsection{Coût total}

Le coût d'EX9 ne peut pas être évalué uniquement par le nombre
d'opérations arithmétiques.

On définit :

\[
\begin{aligned}
T_{\mathrm{total}}
={}&T_{\mathrm{compute}}
+T_{\mathrm{coord}}
+T_{\mathrm{compose}}\\
&+T_{\mathrm{memory}}
+T_{\mathrm{communication}}.
\end{aligned}
\]

Pour une réalisation matérielle, on ajoute notamment :

\[
E_{\mathrm{total}}
=
E_{\mathrm{compute}}
+E_{\mathrm{memory}}
+E_{\mathrm{communication}}.
\]

La réduction du calcul arithmétique peut être annulée par le coût
des transferts, de la coordination ou de la mémoire.

\subsection{Réutilisation}

Soit $N_{\mathrm{naive}}$ le nombre d'exécutions d'opérations dans
une version sans réutilisation, et $N_{\mathrm{reuse}}$ le nombre
d'exécutions après partage des résultats.

On définit :

\[
R_{\mathrm{reuse}}
=
1-
\frac{N_{\mathrm{reuse}}}{N_{\mathrm{naive}}}.
\]

Cette mesure ne suffit pas à démontrer un gain de latence :
les opérations n'ont pas toutes le même coût.

Une mesure pondérée est :

\[
C_{\mathrm{ops}}
=
\sum_{i=1}^{N}
c_i n_i,
\]

où $c_i$ est le coût estimé de l'opération $i$ et $n_i$ son nombre
d'exécutions.

\subsection{Hypothèses falsifiables}

\begin{description}
  \item[H1 -- Réutilisation :]
  le stockage des intermédiaires réduit le nombre de sous-calculs
  répétés sur des graphes comportant des dépendances partagées.

  \item[H2 -- Spécialisation :]
  des sondes spécialisées peuvent réaliser certaines transformations
  avec un coût inférieur à une unité généraliste de référence.

  \item[H3 -- Composition :]
  une composition explicite peut améliorer l'organisation et
  l'ordonnancement des opérations.

  \item[H4 -- Coordination :]
  une politique de coordination peut éviter certaines exécutions
  inutiles lorsque les tâches sont conditionnelles.

  \item[H5 -- Core neuronal :]
  un ensemble de sondes adaptatives peut remplacer certains blocs
  MLP à qualité comparable.

  \item[H6 -- Matériel :]
  une correspondance physique entre cores et ressources matérielles
  peut améliorer le compromis latence-énergie sur des charges adaptées.
\end{description}

Chaque hypothèse doit être testée séparément.

% =========================================================
\section{Plan expérimental EX9}
% =========================================================

\subsection{Étape 1 : graphe de calcul logiciel}

Construire un moteur de graphe où :
\begin{itemize}
  \item chaque opération est un noeud ;
  \item chaque dépendance est une arête ;
  \item les résultats intermédiaires sont identifiables ;
  \item les opérations indépendantes peuvent être planifiées en parallèle.
\end{itemize}

Comparer :
\begin{enumerate}
  \item calcul direct ;
  \item graphe sans réutilisation ;
  \item graphe avec réutilisation ;
  \item graphe avec exécution conditionnelle.
\end{enumerate}

\subsection{Étape 2 : cores logiques}

Créer quatre cores logiques :
\begin{enumerate}
  \item core de calcul ;
  \item core de composition ;
  \item core de coordination ;
  \item core de stockage temporaire.
\end{enumerate}

Le core de communication peut être ajouté comme cinquième module
si son coût ou sa complexité justifie une séparation.

\subsection{Étape 3 : core neuronal}

Remplacer un bloc MLP par un core neuronal de sondes adaptatives.

Comparer :
\begin{itemize}
  \item MLP de référence ;
  \item core neuronal à sondes RBF ;
  \item core neuronal avec fusion pondérée ;
  \item core neuronal avec coordination adaptative.
\end{itemize}

Les budgets de paramètres, les données, les graines et les conditions
d'évaluation doivent être contrôlés.

\subsection{Étape 4 : coordination dynamique}

Introduire une politique qui peut sélectionner :
\begin{itemize}
  \item les cores à activer ;
  \item les opérations à exécuter ;
  \item les résultats à réutiliser ;
  \item les chemins de communication.
\end{itemize}

Comparer cette politique à une exécution statique de tous les cores.

\subsection{Étape 5 : simulation physique}

Avant de concevoir une puce, simuler :
\begin{itemize}
  \item les unités de calcul ;
  \item les mémoires locales ;
  \item les transferts ;
  \item les délais de communication ;
  \item les ressources parallèles.
\end{itemize}

Une simulation de coût matériel doit inclure les transferts et
la synchronisation, pas uniquement les opérations des sondes.

\subsection{Étape 6 : étude matérielle}

Si les résultats logiciels et simulés sont encourageants, étudier
une implémentation matérielle.

Les paramètres à mesurer comprennent :
\begin{itemize}
  \item surface logique ;
  \item fréquence ;
  \item latence ;
  \item énergie par opération ;
  \item bande passante mémoire ;
  \item coût de communication inter-cores.
\end{itemize}

% =========================================================
\section{Protocole de comparaison}
% =========================================================

\subsection{Baselines}

Les comparaisons doivent inclure :

\begin{enumerate}
  \item MLP classique ;
  \item EX7 V003 ;
  \item graphe de calcul conventionnel ;
  \item graphe avec mémorisation ;
  \item EX9 à cores logiques ;
  \item EX9 avec core neuronal ;
  \item EX9 simulé comme architecture physique.
\end{enumerate}

\subsection{Mesures}

\begin{table}[H]
\centering
\small
\begin{tabularx}{\textwidth}{p{3.0cm}X}
\toprule
Mesure & Objectif \\
\midrule
Accuracy / F1 / log-loss
& Qualité prédictive selon la tâche \\

Nombre de paramètres
& Taille des modèles neuronaux \\

Latence p50, p95, p99
& Coût d'exécution réel \\

Nombre d'opérations
& Quantité de calcul effectuée \\

Taux de réutilisation
& Part des résultats intermédiaires réutilisés \\

Mémoire maximale
& Coût des états et intermédiaires \\

Coût de coordination
& Temps consommé par les décisions et l'ordonnancement \\

Énergie
& Coût énergétique par inférence, si mesurable \\

Robustesse
& Stabilité sur des entrées et distributions variées \\
\bottomrule
\end{tabularx}
\caption{Mesures nécessaires à l'évaluation d'EX9.}
\end{table}

\subsection{Règles expérimentales}

Pour garantir une comparaison valide :
\begin{itemize}
  \item utiliser les mêmes données et partitions ;
  \item conserver les mêmes conditions matérielles ;
  \item rapporter plusieurs graines ;
  \item mesurer l'exécution complète ;
  \item inclure le coût de coordination ;
  \item inclure les transferts mémoire ;
  \item distinguer temps d'entraînement et temps d'inférence ;
  \item rapporter moyenne et dispersion des mesures.
\end{itemize}

% =========================================================
\section{Ablations indispensables}
% =========================================================

Les ablations permettront d'identifier la contribution réelle de
chaque mécanisme.

\begin{longtable}{p{3.2cm}p{4.0cm}p{6.0cm}}
\toprule
Variante & Modification & Question testée \\
\midrule
EX9-A & Sans mémoire temporaire &
La réutilisation réduit-elle le coût ? \\

EX9-B & Sans coordination dynamique &
La sélection conditionnelle apporte-t-elle un gain ? \\

EX9-C & Sans composition hiérarchique &
La composition explicite est-elle utile ? \\

EX9-D & Un seul core généraliste &
La spécialisation apporte-t-elle un avantage ? \\

EX9-E & Core neuronal remplacé par MLP &
Le remplacement est-il compétitif ? \\

EX9-F & Communication non optimisée &
Quel est le coût des transferts ? \\

EX9-G & Exécution statique &
Le routage dynamique est-il utile ? \\
\bottomrule
\end{longtable}

% =========================================================
\section{Risques et limites}
% =========================================================

\subsection{Coût de coordination}

Un système modulaire peut être plus lent qu'un bloc dense si la
coordination, les appels de fonctions ou les transferts dominent
le coût total.

\subsection{Granularité}

Des sondes trop petites peuvent créer un excès de communication.
Des sondes trop grandes peuvent perdre l'avantage de spécialisation.

La granularité optimale doit être recherchée expérimentalement.

\subsection{Mémoire temporaire}

La mémorisation peut augmenter l'empreinte mémoire et introduire
des coûts de recherche, d'invalidation et de synchronisation.

\subsection{Approximation neuronale}

Un core neuronal peut produire une approximation différente du MLP.
Il faut mesurer la qualité, la stabilité et la généralisation.

\subsection{Difficulté du passage au matériel}

Une architecture logique efficace ne devient pas automatiquement
une architecture physique efficace. Les contraintes de placement,
de routage, de mémoire et d'énergie peuvent modifier complètement
le résultat.

% =========================================================
\section{Positionnement scientifique}
% =========================================================

EX9 combine plusieurs idées déjà connues dans l'informatique :
\begin{itemize}
  \item graphes de calcul et graphes de dépendances ;
  \item unités de calcul spécialisées ;
  \item mémoire locale et réutilisation des résultats ;
  \item ordonnancement des tâches ;
  \item architectures parallèles ;
  \item modèles neuronaux modulaires ;
  \item approximation de blocs neuronaux.
\end{itemize}

La nouveauté éventuelle ne peut donc pas être déduite de la seule
présence de cores ou de sondes.

Elle devra être établie par :
\begin{enumerate}
  \item une définition précise de l'organisation des sondes ;
  \item une règle originale de composition ou de coordination ;
  \item une comparaison avec les architectures existantes ;
  \item des résultats reproductibles ;
  \item des ablations démontrant la contribution de chaque mécanisme.
\end{enumerate}

La contribution potentielle d'EX9 réside dans l'étude d'une
architecture unifiée où les sondes constituent les éléments
fonctionnels internes de cores logiques, physiques ou neuronaux.

% =========================================================
\section{Feuille de route}
% =========================================================

\begin{table}[H]
\centering
\small
\begin{tabularx}{\textwidth}{p{2.0cm}p{4.0cm}X}
\toprule
Phase & Livrable & Critère de validation \\
\midrule
EX9.1
& Moteur de graphe de calcul
& Exactitude fonctionnelle et dépendances correctes \\

EX9.2
& Cores logiques
& Mesure des coûts de composition et coordination \\

EX9.3
& Mémoire temporaire
& Réutilisation correcte et gain mesuré \\

EX9.4
& Core neuronal
& Comparaison équitable avec MLP \\

EX9.5
& Coordination adaptative
& Réduction mesurée du travail exécuté \\

EX9.6
& Simulation physique
& Modèle de latence, mémoire et communication \\

EX9.7
& Étude matérielle
& Faisabilité et compromis énergie-performance \\
\bottomrule
\end{tabularx}
\caption{Plan de développement proposé pour EX9.}
\end{table}

% =========================================================
\section{Conclusion}
% =========================================================

EX9 propose de considérer le calcul comme un système organisé
d'opérations spécialisées et composables.

Les sondes constituent les éléments internes des cores.
Les cores sont spécialisés selon leur fonction :
calcul, composition, coordination, stockage ou traitement neuronal.

Trois réalisations sont distinguées :
\begin{itemize}
  \item le core logique, qui organise les sondes dans un graphe logiciel ;
  \item le core physique, qui associe les sondes à des ressources matérielles ;
  \item le core neuronal, qui utilise des sondes adaptatives pour remplacer
        des blocs MLP.
\end{itemize}

L'hypothèse centrale est que cette organisation pourrait réduire
les calculs redondants, favoriser la réutilisation et permettre
une spécialisation plus fine.

Cependant, la modularité ne garantit ni une accélération ni une
réduction énergétique. Les coûts de coordination, de mémoire et
de communication doivent être intégrés à chaque mesure.

EX9 doit donc être développé progressivement : d'abord sous forme
de graphe logiciel vérifiable, ensuite sous forme de cores logiques,
puis de cores neuronaux et enfin de simulation et d'étude matérielles.

% =========================================================
\begin{thebibliography}{99}

\bibitem{ex7}
Projet EX,
\textit{Architecture V003 -- OuterCoord-Lite-7x1},
rapport technique interne, archive
\texttt{sondes\_ex7\_v003}.

\bibitem{rbf}
D. S. Broomhead and D. Lowe,
\textit{Multivariable Functional Interpolation and Adaptive Networks},
Complex Systems, vol. 2, pp. 321--355, 1988.

\bibitem{moe}
R. A. Jacobs, M. I. Jordan, S. J. Nowlan, and G. E. Hinton,
\textit{Adaptive Mixtures of Local Experts},
Neural Computation, vol. 3, no. 1, pp. 79--87, 1991.

\bibitem{flynn}
M. J. Flynn,
\textit{Some Computer Organizations and Their Effectiveness},
IEEE Transactions on Computers, vol. C-21, no. 9,
pp. 948--960, 1972.

\bibitem{dataflow}
J. B. Dennis,
\textit{First Version of a Data Flow Procedure Language},
Programming Symposium, Lecture Notes in Computer Science,
vol. 19, pp. 362--376, 1974.

\bibitem{roofline}
S. Williams, A. Waterman, and D. Patterson,
\textit{Roofline: An Insightful Visual Performance Model for
Multicore Architectures},
Communications of the ACM, vol. 52, no. 4, pp. 65--76, 2009.

\bibitem{moe_survey}
C. R. B. et al.,
\textit{A Review of Sparse Mixture of Experts},
Référence à compléter par une revue bibliographique vérifiée
avant toute soumission scientifique.

\end{thebibliography}

\end{document}